\documentclass{article}
\usepackage[utf8]{inputenc}
\usepackage{amsmath,amssymb}
\usepackage[most]{tcolorbox}
\usepackage{geometry}
\geometry{margin=2.5cm}

\newcommand{\step}[1]{\par\noindent\hangindent=3.5em\hangafter=1%
  \makebox[3.5em][l]{\textbf{#1.}}}
\newcommand{\steptag}[1]{\hfill{\small\textsf{[#1]}}}

\newtcolorbox{proofbox}[1]{
  colback=gray!5, colframe=gray!60,
  fonttitle=\bfseries, title={#1},
  breakable, enhanced,
  left=4pt, right=4pt, top=4pt, bottom=4pt,
  before skip=10pt, after skip=10pt
}

\begin{document}

\begin{proofbox}{Compatible Stage-1 polar witnesses and range: there exist...}
\step{1} Compatible Stage-1 polar witnesses and range: there exists one universal def-stage1-polar-witness-data tuple W = (C\_rect, C\_ch, C\_pol, C\_grp, C\_path, C\_der, e\_rect, kappa\_ch, kappa\_pol, kappa\_der, delta\_*, epsilon\_*\textasciicircum{}r, e\_S1, r\_iso), with C\_rect, C\_ch, C\_pol, C\_grp, C\_path, C\_der >= 1, 0 < e\_rect <= 1/C\_rect, and 0 < kappa\_ch, kappa\_pol, kappa\_der <= 1/2, such that all of the following hold simultaneously: (A\_1) for every finite-dimensional epsilon\_X-C*-algebra (calX, I\_X, ., dagger) with 0 <= epsilon\_X <= e\_rect, there are on the same involutive normed space a bilinear product bold-dot and an element J = J\textasciicircum{}dagger for which (calX, J, bold-dot, dagger) satisfies every exact-unit epsilon\_r-C*-algebra axiom of def-epsilon-cstar-algebra, including ||J|| = 1, where epsilon\_r = C\_rect*epsilon\_X, and for every x, y in calX, ||J - I\_X|| <= C\_rect*epsilon\_X and ||x bold-dot y - xy|| <= C\_rect*epsilon\_X*||x||*||y||; (A\_2) for every finite-dimensional exact-unit epsilon\_r-C*-algebra (calX, J, bold-dot, dagger), every delta > 0 with C\_ch*(epsilon\_r + delta) <= kappa\_ch, every V in calUbar\_delta, and every A\textasciicircum{}par in B\textasciicircum{}\{icalH\}\_\{2delta\}(0), there is a unique g\_V(A\textasciicircum{}par) in B\textasciicircum{}\{calH\}\_\{2delta\}(0) such that f\_V(A\textasciicircum{}par + g\_V(A\textasciicircum{}par)) = 0, where f\_V(A) = (1/2)*(((J + A\textasciicircum{}dagger) bold-dot V\textasciicircum{}dagger) bold-dot (V bold-dot (J + A)) - J), the element V bold-dot (J + A\textasciicircum{}par + g\_V(A\textasciicircum{}par)) lies in calU, ||g\_V(A\textasciicircum{}par) + (1/2)*(V\textasciicircum{}dagger bold-dot V - J)|| <= C\_ch*(epsilon\_r*delta + delta\textasciicircum{}2), ||Dg\_V(A\textasciicircum{}par)|| <= C\_ch*(epsilon\_r + delta), and ||D\_\{A\textasciicircum{}perp\} f\_V(A\textasciicircum{}par + g\_V(A\textasciicircum{}par)) - I\_\{calH\}|| <= C\_ch*(epsilon\_r + delta) < 1, and these C\textasciicircum{}1 graph charts cover calU; (A\_3) for every finite-dimensional exact-unit epsilon\_r-C*-algebra, every delta > 0 with C\_ch*(epsilon\_r + delta) <= kappa\_ch, and every family g = (g\_U)\_\{U in calU\} of C\textasciicircum{}1 maps g\_U: B\textasciicircum{}\{icalH\}\_\{2delta\}(0) -> B\textasciicircum{}\{calH\}\_\{2delta\}(0) such that, for every U in calU and A\textasciicircum{}par in B\textasciicircum{}\{icalH\}\_\{2delta\}(0), g\_U(A\textasciicircum{}par) is the unique element of B\textasciicircum{}\{calH\}\_\{2delta\}(0) satisfying f\_U(A\textasciicircum{}par + g\_U(A\textasciicircum{}par)) = 0, where f\_U(A) = (1/2)*(((J + A\textasciicircum{}dagger) bold-dot U\textasciicircum{}dagger) bold-dot (U bold-dot (J + A)) - J), every tangent space T\_U calU is the image of L\_U(I + Dg\_U(0)): icalH -> calX, and omega\_U(Z) = (L\_U\textasciicircum{}\{-1\} Z)\textasciicircum{}par : T\_U calU -> icalH is a global C\textasciicircum{}1 bundle trivialization with distortion at most 1 + C\_ch*epsilon\_r, satisfying omega\_\{cU\}(cZ) = omega\_U(Z) and omega\_U(iU) = iJ for every U in calU, Z in T\_U calU, and c in U(1); (A\_4) for every finite-dimensional exact-unit epsilon\_r-C*-algebra and every delta > 0 with C\_pol*(epsilon\_r + delta) <= kappa\_pol, the map Pi\_delta: calU x B\textasciicircum{}\{calH\}\_delta(J) -> calX, Pi\_delta(U, H) = U bold-dot H, is a C\textasciicircum{}1 diffeomorphism onto the open set S\_delta := Pi\_delta(calU x B\textasciicircum{}\{calH\}\_delta(J)), its inverse (u\_delta, h\_delta): S\_delta -> calU x B\textasciicircum{}\{calH\}\_delta(J) satisfies X = u\_delta(X) bold-dot h\_delta(X), u\_delta(U) = U, and h\_delta(U) = J for every X in S\_delta and U in calU, and calU\_\{delta - C\_pol*(epsilon\_r*delta + delta\textasciicircum{}2)\} subseteq S\_delta subseteq calU\_\{delta + C\_pol*(epsilon\_r*delta + delta\textasciicircum{}2)\}; (A\_5) for every finite-dimensional exact-unit epsilon\_r-C*-algebra and every delta > 0 satisfying C\_pol*(epsilon\_r + delta) <= kappa\_pol and C\_grp*epsilon\_r < delta - C\_pol*(epsilon\_r*delta + delta\textasciicircum{}2), writing (u\_delta, h\_delta) for the unique inverse of Pi\_delta: calU x B\textasciicircum{}\{calH\}\_delta(J) -> S\_delta := Pi\_delta(calU x B\textasciicircum{}\{calH\}\_delta(J)), Pi\_delta(U, H) = U bold-dot H, the formulas mu(U, V) = u\_delta(U bold-dot V) and sigma(U) = u\_delta(U\textasciicircum{}dagger) define C\textasciicircum{}1 maps on all of calU x calU and calU, respectively, and for every U, V, Z in calU, mu(J, U) = mu(U, J) = U, sigma(J) = J, ||mu(U, V) - U bold-dot V|| <= C\_grp*epsilon\_r, ||sigma(U) - U\textasciicircum{}dagger|| <= C\_grp*epsilon\_r, ||mu(mu(U, V), Z) - mu(U, mu(V, Z))|| <= C\_grp*epsilon\_r, ||mu(sigma(U), U) - J|| <= C\_grp*epsilon\_r, and ||mu(U, sigma(U)) - J|| <= C\_grp*epsilon\_r; (A\_6) for every finite-dimensional exact-unit epsilon\_r-C*-algebra, every delta > 0 with C\_pol*(epsilon\_r + delta) <= kappa\_pol, every U\_0, U\_1 in calU, and every q in [0, 1] satisfying ||U\_1 - U\_0|| <= q, C\_path*q <= 1/4, and C\_path*(q + epsilon\_r*q + q\textasciicircum{}2) < delta - C\_pol*(epsilon\_r*delta + delta\textasciicircum{}2), every L\_\{Z\_t\} is invertible and every Z\_t = (1-t)*U\_0 + t*U\_1 lies in calUbar\_\{C\_path*(q + epsilon\_r*q + q\textasciicircum{}2)\} for t in [0, 1], and, writing u\_delta for the unique first component of the inverse of Pi\_delta: calU x B\textasciicircum{}\{calH\}\_delta(J) -> S\_delta := Pi\_delta(calU x B\textasciicircum{}\{calH\}\_delta(J)), Pi\_delta(U, H) = U bold-dot H, the map H(t, U\_0, U\_1) = u\_delta(Z\_t) is jointly continuous in (t, U\_0, U\_1), joins U\_0 to U\_1, and satisfies H(t, cU\_0, cU\_1) = c*H(t, U\_0, U\_1) for every c in U(1); (A\_7) for every finite-dimensional exact-unit epsilon\_r-C*-algebra, every delta > 0, every s in \{+1, -1\}, and every 0 < r <= delta satisfying C\_ch*(epsilon\_r + delta) <= kappa\_ch, C\_pol*(epsilon\_r + delta) <= kappa\_pol, C\_grp*epsilon\_r < delta - C\_pol*(epsilon\_r*delta + delta\textasciicircum{}2), C\_der*(epsilon\_r + r) <= kappa\_der, and (1 + epsilon\_r)*(1 + C\_ch*(epsilon\_r + delta))*r + C\_grp*epsilon\_r < 2*delta, writing u\_delta for the unique first component of the inverse of Pi\_delta: calU x B\textasciicircum{}\{calH\}\_delta(J) -> S\_delta := Pi\_delta(calU x B\textasciicircum{}\{calH\}\_delta(J)), Pi\_delta(U, H) = U bold-dot H, and g\_\{sJ\}: B\_\{2delta\}\textasciicircum{}\{icalH\}(0) -> B\_\{2delta\}\textasciicircum{}\{calH\}(0) for the unique C\textasciicircum{}1 map such that, for every A in B\_\{2delta\}\textasciicircum{}\{icalH\}(0), f\_\{sJ\}(A + g\_\{sJ\}(A)) = 0, where f\_\{sJ\}(B) = (1/2)*(((J + B\textasciicircum{}dagger) bold-dot (sJ)\textasciicircum{}dagger) bold-dot (sJ bold-dot (J + B)) - J), define chi\_s(A) = sJ bold-dot (J + A + g\_\{sJ\}(A)) and the global C\textasciicircum{}1 map sigma(U) = u\_delta(U\textasciicircum{}dagger); then sigma maps chi\_s(B\_r\textasciicircum{}\{icalH\}(0)) into the same sJ-graph chart and, with F\_s(A) = phi\_\{sJ\}\textasciicircum{}par(sigma(chi\_s(A))), one has ||D(F\_s - id)(A) + 2*I\_\{icalH\}|| <= C\_der*(epsilon\_r + r) for every A in B\_r\textasciicircum{}\{icalH\}(0); and (R) delta\_* = min\{1/4, kappa\_ch/(4*C\_ch), kappa\_pol/(4*C\_pol)\}, epsilon\_*\textasciicircum{}r = min\{1/4, kappa\_ch/(4*C\_ch), kappa\_pol/(4*C\_pol), kappa\_der/(8*C\_der), 1/C\_grp, delta\_*/(12*C\_path*C\_grp)\}, e\_S1 = min\{e\_rect, epsilon\_*\textasciicircum{}r/C\_rect\}, r\_iso = min\{delta\_*/4, kappa\_der/(8*C\_der)\}, and for every 0 <= epsilon\_X <= e\_S1, on setting epsilon\_r = C\_rect*epsilon\_X, q = C\_grp*epsilon\_r, r\_- = delta\_* - C\_pol*(epsilon\_r*delta\_* + delta\_*\textasciicircum{}2), and eta = C\_path*(q + epsilon\_r*q + q\textasciicircum{}2), one has C\_ch*(epsilon\_r + delta\_*) <= kappa\_ch, C\_pol*(epsilon\_r + delta\_*) <= kappa\_pol, q < r\_-, C\_path*q <= 1/4, eta < r\_-, C\_der*(epsilon\_r + r\_iso) <= kappa\_der, (1 + epsilon\_r)*(1 + C\_ch*(epsilon\_r + delta\_*))*r\_iso + q < 2*delta\_*, r\_- >= 3*delta\_*/4, eta <= delta\_*/4, and C\_der*(r\_iso + epsilon\_r) <= kappa\_der/4 < 1. \steptag{validated} \\[6pt]
\step{1.1} Final conjunction assembly: choose the single common Stage-1 polar witness tuple supplied by the tuple-selection child; use the seven transport children to obtain clauses (A\_1)-(A\_7) for that same tuple and the scalar-arithmetic child to obtain clause (R), and conjoin the tuple range, (A\_1)-(A\_7), and (R) to prove the root contract. \steptag{validated} \\[6pt]
\step{1.1.1} Common tuple selection and range: take the existential base coefficients and positive margins supplied by lem-stage1-rectified-cstar-transport, lem-stage1-unitary-graph-transport, lem-stage1-maurer-cartan-transport, lem-stage1-polar-retraction-transport, lem-stage1-approximate-group-laws-transport, lem-stage1-polar-path-transport, and lem-stage1-inversion-derivative-transport. Define each of C\_rect,C\_ch,C\_pol,C\_grp,C\_path,C\_der as the maximum of 1 and every base coefficient bearing that name; define e\_rect as the minimum of the rectification margin and 1/C\_rect; define each kappa\_ch,kappa\_pol,kappa\_der as the minimum of 1/2 and every positive base margin bearing that name. These finite maxima and minima give coefficients >=1, 0<e\_rect<=1/C\_rect, 0<kappa\_ch,kappa\_pol,kappa\_der<=1/2, and meet every helper's coefficient lower bounds and margin upper bounds. Define delta\_*, epsilon\_*\textasciicircum{}r, e\_S1, and r\_iso by the four formulas in clause (R); by def-stage1-polar-witness-data these fourteen scalars form one tuple W. \steptag{validated} \\[4pt]
\step{1.1.2} Rectification application: for the common W satisfying the rectification coefficient and margin bounds from the tuple-selection child, lem-stage1-rectified-cstar-transport applies verbatim and yields clause (A\_1) of the root contract, with epsilon\_r=C\_rect*epsilon\_X, for every finite-dimensional epsilon\_X-C*-algebra and every 0<=epsilon\_X<=e\_rect. \steptag{validated} \\[4pt]
\step{1.1.3} Unitary-graph application: for the common W satisfying the graph coefficient and margin bounds from the tuple-selection child, lem-stage1-unitary-graph-transport applies verbatim and yields clause (A\_2) of the root contract for every finite-dimensional exact-unit epsilon\_r-C*-algebra, every admissible delta,V,A\textasciicircum{}par, including uniqueness, membership, all three displayed estimates, and chart coverage. \steptag{validated} \\[4pt]
\step{1.1.4} Maurer-Cartan application: for the common W satisfying the chart coefficient and margin bounds from the tuple-selection child, lem-stage1-maurer-cartan-transport applies verbatim and yields clause (A\_3) of the root contract for every finite-dimensional exact-unit epsilon\_r-C*-algebra, every admissible delta, and every stipulated family g, including the tangent-space formula, global C\textasciicircum{}1 trivialization, distortion, scalar equivariance, and omega\_U(iU)=iJ. \steptag{validated} \\[4pt]
\step{1.1.5} Polar-retraction application: for the common W satisfying the polar coefficient and margin bounds from the tuple-selection child, lem-stage1-polar-retraction-transport applies verbatim and yields clause (A\_4) of the root contract for every finite-dimensional exact-unit epsilon\_r-C*-algebra and every admissible delta, including the C\textasciicircum{}1 diffeomorphism, inverse identities, and both set inclusions. \steptag{validated} \\[4pt]
\step{1.1.6} Approximate-group application: for the common W satisfying the group and polar coefficient and margin bounds from the tuple-selection child, lem-stage1-approximate-group-laws-transport applies verbatim and yields clause (A\_5) of the root contract for every finite-dimensional exact-unit epsilon\_r-C*-algebra and every delta satisfying its two displayed hypotheses, including global C\textasciicircum{}1 definitions of mu,sigma, exact unit identities, and all five C\_grp*epsilon\_r estimates. \steptag{validated} \\[4pt]
\step{1.1.7} Polar-path application: for the common W satisfying the path and polar coefficient and margin bounds from the tuple-selection child, lem-stage1-polar-path-transport applies verbatim and yields clause (A\_6) of the root contract for every finite-dimensional exact-unit epsilon\_r-C*-algebra and all delta,U\_0,U\_1,q satisfying its displayed hypotheses, including invertibility, approximate-unitarity of Z\_t, and the jointly continuous equivariant joining path. \steptag{validated} \\[4pt]
\step{1.1.8} Inversion-derivative application: for the common W satisfying the derivative, chart, polar, and group coefficient and margin bounds from the tuple-selection child, lem-stage1-inversion-derivative-transport applies verbatim and yields clause (A\_7) of the root contract for every finite-dimensional exact-unit epsilon\_r-C*-algebra and all delta,s,r satisfying its displayed hypotheses, including preservation of the sJ graph chart and the stated derivative estimate for F\_s. \steptag{validated} \\[4pt]
\step{1.1.9} Scalar-arithmetic application: the common tuple has coefficients >=1, 0<e\_rect<=1/C\_rect, and 0<kappa\_ch,kappa\_pol,kappa\_der<=1/2, and its four derived scales were defined by the displayed formulas; therefore lem-stage1-polar-scalar-arithmetic applies verbatim and yields every assertion of clause (R), with epsilon\_r,q,r\_-,eta defined there, for every 0<=epsilon\_X<=e\_S1. \steptag{validated} \\[4pt]
\end{proofbox}

\end{document}
